\chapter{Direct page retracing}\label{module:retrace}
\Q{Evidence}{Current implementation boundary}
\path{Software/DemoCode/retrace.py}. Exact file hashes are in the accompanying source manifest. Names, declarations and source locations below come from the captured source, not a runtime import.
\Q{Purpose}{Input-to-output reading flow}
Take page pixels through line segmentation, skeleton tracing and pruning, then map recovered trajectories for reproduction. The current SegmentLines call mismatch is described in the pipeline audit; this is not a verified successful run.


\Q{Evidence}{Selected current module settings}
These are declarations or expressions, not recovered environment values. Source comments do not establish a measured optimum. Larger glyph coordinate tables are designed data in the bundled source.
\begin{lstlisting}[style=plain]
HERE = Path(__file__).resolve().parent
MIN_STROKE_PTS = 4
MIN_STROKE_XH = 0.12
\end{lstlisting}
\section*{Stage: Direct ink retracing}
\section{Method background: Direct ink retracing}\label{method:retrace}
\Q{Purpose}{The role in this system}
Bypass recognition and glyph synthesis to compare a traced source path with physical output.


\Q{Justification}{The reason for this construction}
It removes two sources of variation but still includes binarisation, thinning, tracing, pruning, fitting and motion; it is not a pure mechanical ground truth.


\Q{Alternatives}{What changes if another method is used}
Direct known geometric paths isolate motion more tightly. OCR-and-synthesise asks a different end-to-end question. An online digitiser would supply a better observed trajectory.


\Q{Library}{Implementation speed and accuracy}
Image/vector tracing libraries could replace centreline extraction, with different topology and simplification. The current PageStrokes timer keyword is incompatible with current SegmentLines and is documented as a static integration discrepancy. \cite{code2026,zhangsuen1984,douglaspeucker1973}.


\section{PageStrokes}\label{fn:retrace:PageStrokes}
\begin{lstlisting}[style=sig]
def PageStrokes(gray, timer=None)
\end{lstlisting}
\Q{Evidence}{Source and method link}
\path{Software/DemoCode/retrace.py} line 49. Implementation evidence: \cite{code2026}. Related method, alternatives and library discussion: Section~\ref{method:retrace}.


\Q{Purpose}{What this contributes}
Calls SegmentLines(gray, timer=timer), although the captured core does not accept timer. If that mismatch were resolved elsewhere, the remaining path would still binarise, thin, trace and prune source ink; the code comment claiming only mechanical differences is stronger than this processing supports.


\Q{Method}{Current implementation}
Bypass recognition and glyph synthesis to compare a traced source path with physical output. The excerpt below is from the current definition; the complete function is in the source snapshot. Its constants and branches, rather than an older signature, govern this path.


\begin{lstlisting}[style=plain]
def PageStrokes(gray, timer=None):
    lines, _dbg, info = SEG.SegmentLines(gray, timer=timer)
    strokes = []
    perLine = []
    for ln in lines:
        crop = np.asarray(ln["image"], dtype=np.uint8)
        ink = SEG.BinarizeLine(crop)
        if not ink.any():
            perLine.append(0)
            continue
        polys = ST.TracePolylines(SEG.Skeletonize(ink))
        x0, y0 = (ln["bbox"][0], ln["bbox"][1]) if "bbox" in ln else (0, 0)
        kept = 0
        for p in polys:
            if len(p) < MIN_STROKE_PTS:
                continue
            strokes.append([(float(x)+x0, -(float(y)+y0)) for (x, y) in p])
            kept += 1
        perLine.append(kept)
    return strokes, dict(lines=len(lines), strokesPerLine=perLine, **info)
\end{lstlisting}
\Q{Justification}{Interpret the method rather than infer a benchmark}
It removes two sources of variation but still includes binarisation, thinning, tracing, pruning, fitting and motion; it is not a pure mechanical ground truth. \cite{code2026,zhangsuen1984,douglaspeucker1973}.


\section{LinePitchMm}\label{fn:retrace:LinePitchMm}
\begin{lstlisting}[style=sig]
def LinePitchMm(info, cfg, traj)
\end{lstlisting}
\Q{Evidence}{Source and method link}
\path{Software/DemoCode/retrace.py} line 77. Implementation evidence: \cite{code2026}. Related method, alternatives and library discussion: Section~\ref{method:retrace}.


\Q{Purpose}{What this contributes}
The segmenter reports the pitch in half-resolution pixels; the trajectory is in full-resolution page pixels, so the pitch doubles, and ToMachine's own scale converts it to millimetres.


\Q{Method}{Current implementation}
Bypass recognition and glyph synthesis to compare a traced source path with physical output. The excerpt below is from the current definition; the complete function is in the source snapshot. Its constants and branches, rather than an older signature, govern this path.


\begin{lstlisting}[style=plain]
def LinePitchMm(info,cfg,traj):
    _p,scale = REP.ToMachine(traj,cfg,optimizeOrder=False)
    return max(1.0,float(info.get("pitch",40))*2.0*scale)
\end{lstlisting}
\Q{Justification}{Interpret the method rather than infer a benchmark}
It removes two sources of variation but still includes binarisation, thinning, tracing, pruning, fitting and motion; it is not a pure mechanical ground truth. \cite{code2026,zhangsuen1984,douglaspeucker1973}.


\section{PruneSpecks}\label{fn:retrace:PruneSpecks}
\begin{lstlisting}[style=sig]
def PruneSpecks(strokes, xhPx)
\end{lstlisting}
\Q{Evidence}{Source and method link}
\path{Software/DemoCode/retrace.py} line 83. Implementation evidence: \cite{code2026}. Related method, alternatives and library discussion: Section~\ref{method:retrace}.


\Q{Purpose}{What this contributes}
Bypass recognition and glyph synthesis to compare a traced source path with physical output.


\Q{Method}{Current implementation}
Bypass recognition and glyph synthesis to compare a traced source path with physical output. The excerpt below is from the current definition; the complete function is in the source snapshot. Its constants and branches, rather than an older signature, govern this path.


\begin{lstlisting}[style=plain]
def PruneSpecks(strokes, xhPx):
    limit = MIN_STROKE_XH*max(1.0, xhPx)
    out = []
    for s in strokes:
        span = max(max(p[0] for p in s)-min(p[0] for p in s),
                   max(p[1] for p in s)-min(p[1] for p in s))
        if len(s) >= MIN_STROKE_PTS and span >= limit:
            out.append(s)
    return out
\end{lstlisting}
\Q{Justification}{Interpret the method rather than infer a benchmark}
It removes two sources of variation but still includes binarisation, thinning, tracing, pruning, fitting and motion; it is not a pure mechanical ground truth. \cite{code2026,zhangsuen1984,douglaspeucker1973}.


\section{RetracePage}\label{fn:retrace:RetracePage}
\begin{lstlisting}[style=sig]
def RetracePage(imagePath, outPath, cfg=None, previewPath=None)
\end{lstlisting}
\Q{Evidence}{Source and method link}
\path{Software/DemoCode/retrace.py} line 100. Implementation evidence: \cite{code2026}. Related method, alternatives and library discussion: Section~\ref{method:retrace}.


\Q{Purpose}{What this contributes}
The whole job, photograph in and a g-code file out the trajectory is handed to the ordinary writer, so the page is placed, scaled to the work area and ordered by exactly the same code a synthesised line goes through.  only the source of the strokes differs.


\Q{Method}{Current implementation}
Bypass recognition and glyph synthesis to compare a traced source path with physical output. The excerpt below is from the current definition; the complete function is in the source snapshot. Its constants and branches, rather than an older signature, govern this path.


\begin{lstlisting}[style=plain]
def RetracePage(imagePath, outPath, cfg=None, previewPath=None):
    cfg = cfg or REP.GantryConfig()
    run = TIMING.Run("retrace")

    with run.Step("loading"):
        gray = TR.ReadImageFile(imagePath)

    strokes, info = PageStrokes(gray, timer=run)
    with run.Step("tracing"):
        strokes = PruneSpecks(strokes, info.get("pitch", 40)*0.4)
    if not strokes:
        raise SystemExit("no ink found to trace in %s" % imagePath)

    traj = SYN.Trajectory(strokes, dict(author="retrace", source=str(imagePath),
                                        lines=info["lines"]))
    with run.Step("gcode"):
        written = REP.WriteGcode(traj, cfg, str(outPath),
                                 title="retrace of %s" % Path(imagePath).name,
                                 rowTolMm=LinePitchMm(info, cfg, traj))

    preview = None
    if previewPath:
        with run.Step("preview"):
            preview = SYN.RenderTrajectory(traj, pxPerMm=6.0)
            from PIL import Image
            Image.fromarray(np.asarray(preview, dtype=np.uint8)).save(previewPath)

    run.chars = 0
    run.lines = info["lines"]
    out = dict(written)
    out.update(gcode=str(outPath), tracedStrokes=len(strokes), timing=run.AsDict(),
               preview=str(previewPath) if previewPath else None)
    return out
\end{lstlisting}
\Q{Justification}{Interpret the method rather than infer a benchmark}
It removes two sources of variation but still includes binarisation, thinning, tracing, pruning, fitting and motion; it is not a pure mechanical ground truth. \cite{code2026,zhangsuen1984,douglaspeucker1973}.


\section{Main}\label{fn:retrace:Main}
\begin{lstlisting}[style=sig]
def Main()
\end{lstlisting}
\Q{Evidence}{Source and method link}
\path{Software/DemoCode/retrace.py} line 141. Implementation evidence: \cite{code2026}. Related method, alternatives and library discussion: Section~\ref{method:retrace}.


\Q{Purpose}{What this contributes}
Bypass recognition and glyph synthesis to compare a traced source path with physical output.


\Q{Method}{Current implementation}
Bypass recognition and glyph synthesis to compare a traced source path with physical output. The excerpt below is from the current definition; the complete function is in the source snapshot. Its constants and branches, rather than an older signature, govern this path.


\begin{lstlisting}[style=plain]
def Main():
    ap = argparse.ArgumentParser(description=__doc__.splitlines()[0])
    ap.add_argument("image")
    ap.add_argument("-o", "--out", default=None)
    ap.add_argument("--preview", default=None)
    args = ap.parse_args()
    out = args.out or str(Path(args.image).with_suffix(".gcode").name)
    r = RetracePage(args.image, out, previewPath=args.preview)
    print("traced %d lines into %d strokes" % (r["timing"]["lines"], r["tracedStrokes"]))
    print("  %s: %d g-code lines, %d pen pulses, scale %.3f"
          % (r["gcode"], r["lines"], r["penPulses"], r["scale"]))
    if r["clipped"]:
        print("  WARNING: %d point(s) clamped to the work area" % r["clipped"])
    if r["preview"]:
        print("  preview: %s" % r["preview"])
\end{lstlisting}
\Q{Justification}{Interpret the method rather than infer a benchmark}
It removes two sources of variation but still includes binarisation, thinning, tracing, pruning, fitting and motion; it is not a pure mechanical ground truth. \cite{code2026,zhangsuen1984,douglaspeucker1973}.

